% Replacement for inherited A.CR2 / paper A.15.
% Checked by GPT-6 Astra Pro, 2026-11-18.
% Pro mode; separate effort setting not exposed.
\subsection{Compactification and point counts for an arbitrary retained prime set}

\begin{claim}{A.15}{Proved}{The full compactified pro-$r$ group at every open level}
Let $C/k$ be a smooth proper geometrically connected curve over a
finite field of characteristic $p$, let $P\in C(k)$, and let
$r\in\Sigma\setminus\{p\}$, with no other restriction on $\Sigma$.
Put $U=C-P$, $F=\pi_1(U_{\bar k})^\Sigma$,
$\Delta=\pi_1(C_{\bar k})^\Sigma$, $K=\ker(F\to\Delta)$, and
$N=\ker(K\to K^{(r)})$. Let $G_P$ be the arithmetic relative
cuspidally pro-$r$ quotient, with geometric kernel $F/N$ and full
arithmetic quotient $G_k$.

For every open geometric subgroup $I/N\le F/N$, let $V$ be its
actual connected geometric cover of $U_{\bar k}$ and let $\bar V$
be its smooth compactification. If $T_I$ is the closed normal
subgroup of $I/N$ generated by all its cusp inertia images, then
the actual covering and compactification maps induce an isomorphism
\[
             ((I/N)/T_I)^{(r)}\simeq\pi_1(\bar V)^{(r)}.
\]
This isomorphism is compatible with every arithmetic outer action
preserving the cover. Consequently, a cusp-preserving isomorphism
of such arithmetic relative groups over $G_k$ preserves the
integral Frobenius polynomial and the point counts of every pair
of corresponding open covers, over every finite constant extension.
In particular these assertions allow $p\in\Sigma$.
\end{claim}

\begin{proof}
Write $\Gamma=\pi_1(U_{\bar k})$ before taking a prime-set quotient.
We first justify the two geometric facts needed at an open level.

\paragraph{Actual covers, tameness, and the induced $\Sigma$ topology.}
The inverse image $\widetilde I$ of $I$ in $\Gamma$ is open,
and hence defines an actual connected finite \emph{\'etale} cover
$V\to U_{\bar k}$. Its Galois closure has group a finite quotient
of $F/N$. The image of every cusp inertia subgroup lies in the
image of $K/N$, a pro-$r$ group. Its finite inertia groups
therefore have $r$-power order and are tame because $r\ne p$.
Thus the cover extends by normalization to a finite tame morphism
of smooth proper curves, and its boundary is the finite set
$\bar V-V$. The same reasoning applies to arithmetic covers after
a finite extension of constants. This tameness assertion places
no restriction on the $p$-part of their proper deck groups.

The natural map $\widetilde I^\Sigma\to I$ is an isomorphism.
Here is the full normal-core argument, including $p\in\Sigma$.
The subgroup $B=\operatorname{core}_{\Gamma}(\widetilde I)$ is
open normal, and $\Gamma/B$ is a finite $\Sigma$-group, since
$\widetilde I$ is pulled back from the pro-$\Sigma$ group $F$.
Given any open normal $J\triangleleft\widetilde I$ with
$\widetilde I/J$ a finite $\Sigma$-group, put $J'=J\cap B$ and
\[
                 J_0=\bigcap_{g\in\Gamma/B}gJ'g^{-1}.
\]
Then $J'$ is normal in $B$, $B/J'$ is a finite $\Sigma$-group,
and $J_0$ is open normal in $\Gamma$. The group $B/J_0$ embeds
in the product of the conjugates of $B/J'$. Hence both $B/J_0$
and $\Gamma/J_0$ are finite $\Sigma$-groups, and $J_0\subseteq J$.
Every finite $\Sigma$-quotient of $\widetilde I$ is therefore
detected by a finite $\Sigma$-quotient of $\Gamma$. This proves
the asserted induced topology and the isomorphism. It uses only
that finite $\Sigma$-groups are closed under subgroups, finite
products, and extensions.

\paragraph{Killing the actual boundary at the retained prime.}
Let $T\triangleleft I$ be the closed normal subgroup generated by
the images of the actual inertia subgroups of $\widetilde I$ at
the points of $\bar V-V$. A finite \emph{\'etale} cover of $V$
extends \emph{\'etale} over $\bar V$ exactly when each of those
inertia groups acts trivially. Indeed its normalization is finite
over the smooth curve, and the criterion is the unramified
criterion at each discrete valuation ring. The universal property
of pro-$\Sigma$ completion therefore gives
\[
                       I/T\simeq\pi_1(\bar V)^\Sigma.
\]
This is right exactness of completion, and uses no left-exactness
claim for completion.

The same-prime quotient lemma A.14 applies to $N\subset I$:
every continuous homomorphism from $I$ to a finite $r$-group
kills $N$. In particular the finite $r$-quotients of $I$ killing
$T$ are exactly those killing the closed product $TN$.
The image of $T$ in $I/N$ is precisely $T_I$. Therefore
\[
 \begin{aligned}
 ((I/N)/T_I)^{(r)}
   &\simeq (I/TN)^{(r)}\\
   &\simeq (I/T)^{(r)}\\
   &\simeq (\pi_1(\bar V)^\Sigma)^{(r)}
    =\pi_1(\bar V)^{(r)}.
 \end{aligned}
\]
In the final equality we used $r\in\Sigma$. Each map is induced
by an actual geometric quotient and is consequently natural.
There is no assumption that $N\subseteq T$, or that $[F:I]$
is an $r$-power.

\paragraph{Arithmetic compatibility and Frobenius.}
Let $L/k$ be finite and let $H\le(G_P)_L$ be an open subgroup
whose arithmetic projection is the whole $G_L$. Its actual
connected cover $V_H\to U_L$ is geometrically connected: the
image of its arithmetic fundamental group in $G_L$ is all of
$G_L$, so its field of constants is $L$. Its geometric subgroup
is an $I/N$ as above. Arithmetic conjugation by $H$ permutes
the actual boundary inertia groups and preserves every kernel
in the displayed construction. The isomorphism is thus equivariant
for the outer $G_L$-action on the geometric compactification group.
It identifies, in particular,
\[
 H^1(\bar V_{H,\bar L},\mathbf Z_r)
   =\operatorname{Hom}_{\mathrm{cont}}
        (\pi_1(\bar V_{H,\bar L})^{(r)},\mathbf Z_r)
\]
with its actual continuous Frobenius action.

An isomorphism of the original arithmetic relative groups over
$G_k$, preserving the cusp decomposition conjugacy class, identifies
the groups belonging to corresponding open subgroups $H$.
It also identifies their systems of cusp inertia. Their compactified
$r$-adic $H^1$ representations are therefore isomorphic over
$G_L$. For the geometric Frobenius $\operatorname{Fr}_L$, the
smooth proper curve trace formula and integrality theorem give
\[
 P_{\bar V_H}(t)=\det(1-t\operatorname{Fr}_L\mid
       H^1(\bar V_{H,\bar L},\mathbf Q_r))\in\mathbf Z[t],
 \qquad
 \#\bar V_H(L_m)=|L|^m+1-
       \operatorname{Tr}(\operatorname{Fr}_L^m\mid H^1),
\]
where $L_m/L$ is the degree-$m$ extension. These are
\citet[formula (1.5.1) and Theorem (1.6)]{WeilCurves}; their
hypotheses apply to every smooth proper geometrically connected
compactification, including genus zero or one. Equality of the
$r$-adic characteristic polynomials implies equality of their
integer polynomials because $\mathbf Z[t]\hookrightarrow\mathbf Q_r[t]$.
The compactified point counts consequently agree for every $m$.

\paragraph{Recovering the deleted points.}
Choose one actual cusp decomposition subgroup $D_P$ of $(G_P)_L$.
The boundary points of $\bar V_H$ are indexed by the finite double
coset set $H\backslash(G_P)_L/D_P$. For a representative $g$, the
residue degree of the corresponding boundary point is
\[
       [G_L:\operatorname{pr}(H\cap gD_Pg^{-1})].
\]
This follows directly from the decomposition-group action on the
finite fibre of the cover: the stabilizer is the displayed
intersection and its arithmetic projection is the Galois group
of the residue field. It remains valid in the relative quotient,
because the cover factors through that quotient and the actual
decomposition action has exactly the same image. The isomorphism
over $G_L$ identifies the double cosets and these degrees.
A deleted point of degree $d$ contributes $d$ points over $L_m$
if $d\mid m$, and zero otherwise. Subtracting the equal boundary
counts from the compactified counts proves
$\#V_H(L_m)=\#V_{H'}(L_m)$ for all corresponding covers and all $m$.
For an open subgroup with a smaller arithmetic projection, begin
over its actual finite field of constants; this covers every open
level claimed in the statement.
\end{proof}

\paragraph{Dependency and check record.}
This replacement has the single internal group-theoretic dependency
A.14, together with actual finite-cover compactification, local
unramifiedness, and the smooth proper curve trace formula. It does
not use A.16, a characteristic-zero lift, an open homotopy exact
sequence, or SGA~1 XIII 2.9--2.11. In particular A.49 can use
A.14--A.15 with the full stated range of $\Sigma$. Checked in this
writing on 18 November 2026 by GPT-6 Astra Pro, Pro mode; a separate
reasoning-effort setting was not exposed.
